\documentclass[12pt]{article}
\usepackage[margin=0.75in]{geometry}

\newcommand{\duedate}{01/28/2026}
\newcommand{\assignment}{Continual Learning Overview: Summary}

\newcommand{\name}{Aksel Kretsinger-Walters, adk2164}
\newcommand{\email}{adk2164@columbia.edu}

\makeatletter
\def\input@path{{../}{../../}{../../../}}
\makeatother
\input{pset_template_CLMM.tex} %% DO NOT CHANGE THIS LINE

\begin{document}
\psetheader %% DO NOT CHANGE THIS LINE

\section*{Titans: Learning to Memorize at Test Time}

\paragraph{Motivation.}
Transformers achieve strong performance due to attention-based context retrieval, but they are fundamentally limited by
their finite context window and quadratic complexity with respect to sequence length. As a result, they struggle with
very long contexts and cannot persistently store information beyond the active attention window.

The Titans architecture introduces a neural long-term memory module that allows models to \emph{learn to memorize at
test time}. The goal is to combine the strengths of attention (accurate short-term dependency modeling) with
a persistent
memory mechanism that can store information over extremely long time horizons.

\paragraph{Core Idea.}
The paper views sequence models through a memory lens:

\begin{itemize}
		\item Attention acts as \textbf{short-term memory}.
		\item Recurrent memory modules act as \textbf{long-term memory}.
\end{itemize}

Titans integrate both components into a unified architecture where attention handles the local context and a separate
memory module accumulates information over long time spans. This allows the model to scale to contexts far beyond
traditional Transformer limits. \cite{titans}

\paragraph{Neural Long-Term Memory.}
The key contribution is a neural memory module that stores information dynamically during inference.
Unlike standard neural networks whose parameters remain fixed during test time, the memory module
updates itself using the incoming data stream.

The system measures the \emph{surprise} of an input token using gradients of an associative memory objective.
Tokens that violate model expectations are considered more informative and are preferentially stored
in memory.

Memory updates resemble gradient-based learning rules:

\[
		M_t = M_{t-1} - \eta \nabla L(M_{t-1}; x_t)
\]

where $M_t$ is the memory state and $L$ is a memory objective that captures key–value associations.

\paragraph{Memory Management.}
To prevent uncontrolled memory growth, the architecture introduces a decay mechanism similar to weight
decay in optimization. This mechanism balances:

\begin{itemize}
		\item the amount of stored information
		\item the surprise level of incoming data
		\item the total memory capacity
\end{itemize}

This effectively implements a learned forgetting mechanism.

\paragraph{Titan Architecture.}
The full model consists of three interacting components:

\begin{itemize}
		\item \textbf{Core module:} a short-context attention mechanism that processes the current sequence window.
		\item \textbf{Long-term memory module:} a neural memory that accumulates information over long horizons.
		\item \textbf{Persistent memory:} task-specific parameters representing global knowledge about the domain.
\end{itemize}

The paper introduces three variants that integrate the memory module in different ways:
as an additional context source, as a network layer, or as a gated branch.

\paragraph{Results.}
Experiments show that Titans:

\begin{itemize}
		\item outperform modern recurrent sequence models
		\item match or exceed Transformers with equivalent context windows
		\item scale to contexts exceeding two million tokens
\end{itemize}

Performance improvements are observed in language modeling, commonsense reasoning,
time-series forecasting, genomics tasks, and long-context retrieval benchmarks. \cite{titans}

\paragraph{Limitations.}
Despite promising results, several limitations remain:

\begin{itemize}
		\item The architecture introduces additional complexity compared to standard Transformers.
		\item Memory updates may introduce instability during long sequences.
		\item The memory module relies on heuristic design choices such as surprise-based storage.
\end{itemize}

Further work is required to evaluate scaling behavior in large language models.

\paragraph{Takeaway.}
Titans propose a new hybrid architecture combining attention-based short-term reasoning with
a learnable long-term memory that updates during inference. This approach aims to bridge the gap
between Transformers' powerful local modeling and the persistent memory capabilities of recurrent systems.

\end{document}
